\clearpage\fancyhead[L]{\small PALLAS.}\hypertarget{native-388-lower}{}\label{native-388-lower}
\begin{center}\textbf{PALLAS.}\par\vspace{4pt}\hrule\vspace{4pt}\fontsize{10}{12}\selectfont Astronomisches Jahrbuch für 1818. Seite 167\;..\;173. Berlin 1815.\par\vspace{4pt}\hrule\vspace{4pt}\end{center}

\begin{flushright}
Göttingen, den 1.\ Februar 1815.
\end{flushright}

\par\indent Die auch im letzten Jahrbuche vorkommenden Berechnungen für die \emph{Pallas} haben sich wiederum auf das beste bestätigt. Ich selbst habe zwar den Planeten nur Einmal mit dem Kreismikrometer beobachtet

\begin{center}
\fontsize{9}{11}\selectfont\setlength{\tabcolsep}{3pt}\renewcommand{\arraystretch}{1.15}\begin{tabular}{l|ccc|ccc|ccc}
1814 & \multicolumn{3}{c|}{Mittlere Zeit} & \multicolumn{3}{c|}{Gerade Aufsteig.} & \multicolumn{3}{c}{Abweichung} \\
\hline
Sept.\ 16 & $12^{\mathrm u}$ & $1^{\mathrm m}$ & $1^{\mathrm s}$ & $146^\circ$ & $16'$ & $23''4$ & $11^\circ$ & $4'$ & $20''9$ S. \\
\end{tabular}
\end{center}

\noindent allein durch eine schätzbare Reihe von Meridianbeobachtungen von den Herren von \textsc{Lindenau}, \textsc{Bessel} und \textsc{Schumacher} wurde ich in den Stand gesetzt die Opposition gut zu berechnen. Wegen der Details beziehe ich mich auf die Gött.\ gel.\ Anz.\ St.\ 199 [S.\ 379 d.\ B.] und führe hier blos das Endresultat an:

\begin{center}1814 Octob. 25. $12^{\mathrm u}\;33^{\mathrm m}\;22^{\mathrm s}$ Mittlere Zeit in Göttingen\end{center}
\begin{center}\fontsize{10}{12}\selectfont\setlength{\tabcolsep}{3pt}\renewcommand{\arraystretch}{1.15}\begin{tabular}{lcr}
Wahre Länge & $.\;.\;.\;.\;.\;.\;.\;.\;.\;.\;.\;.\;.$ & $31^\circ\;58'\;11''3$\\
Geocentrische Breite, südl. & $.\;.\;.\;.\;.\;.\;.\;.$ & $37\;20\;53.2$\\
\end{tabular}\end{center}

\par\indent Die sämmtlichen 10 bisherigen Oppositionen werden durch meine Theorie bis auf wenige Sekunden dargestellt.

\Needspace{12\baselineskip}\par\indent Hr.\ \textsc{Enke} hat die verdienstliche Arbeit über sich genommen, die Rechnungen für die nächste Erscheinung im Voraus zu machen. Ich setze hier sein Resultat für die XI.\ Opposition her:

\begin{center}1816 April 9. $2^{\mathrm u}\;14^{\mathrm m}\;24^{\mathrm s}$ Mittlere Zeit in Göttingen\end{center}
\begin{center}\fontsize{10}{12}\selectfont\setlength{\tabcolsep}{3pt}\renewcommand{\arraystretch}{1.15}\begin{tabular}{lcr}
Wahre Länge & $.\;.\;.\;.\;.\;.\;.\;.\;.\;.\;.\;.$ & $199^\circ\;34'\;40''5$\\
Geocentrische Breite & $.\;.\;.\;.\;.\;.\;.\;.\;.$ & $28\;6\;16.3\;\text{Nordl.}$\\
\end{tabular}\end{center}

\Needspace{8\baselineskip}\par\indent In der nächsten Opposition wird die \emph{Pallas} heller sein, als sie seit mehreren Jahren gewesen ist. Hier die verglichene Lichtstärke

\begin{center}\fontsize{9}{11}\selectfont
1812\ .\ 0.01666,\quad 1813\ .\ 0.01475,\quad 1814\ .\ 0.05476,\quad 1816\ .\ 0.05997
\end{center}

\par\indent Von der Ephemeride, welche derselbe geschickte junge Astronom berechnet hat, habe ich das Vergnügen Ihnen hier eine Abschrift beizulegen.

\par\indent Für die Beobachtungen des Wintersolstitium war in diesem Jahre das Wetter ungewöhnlich günstig. Ich habe nichts unterlassen, um diesen Beobachtungen alle mögliche Genauigkeit zu geben. Dazu \SourceLeaf{389}gehört, dass ein eigner Schirm blos auf das Objectiv Licht fallen liess, und alle andern Strahlen abhielt. Da bei meinen Beobachtungen früherer Solstitien diese Einrichtung noch nicht angebracht war, so wollte ich denselben kein vollkommenes Vertrauen schenken; inzwischen kann ich doch eigentlich nicht sagen, dass die Beobachtungen vor Anbringung des Schirms eine schlechtere Uebereinstimmung gegeben hätten. Das Resultat in Beziehung auf die Schiefe der Ekliptik ist nun ganz übereinstimmend mit dem des Hrn.\ \textsc{von Zach}, nemlich durch das Wintersolstitium finden sich $10''$ weniger als durch das Sommersolstitium. Ich möchte indessen dieses Phänomen lieber so ausdrücken: Wenn man für die Schiefe der Ekliptik den mittlern Werth zum Grunde legt, so geben Sonnenbeobachtungen immer eine im Durchschnitt $5''$ kleinere Polhöhe als Circumpolarsterne, denn nach meinen Erfahrungen gilt dies zu \emph{jeder} Jahrszeit. Dies sonderbare Phaenomen ist noch ein astronomisches Räthsel. Läge in dem Instrument irgend ein Grund, der uns berechtigte anzunehmen, dass es alle Zenithdistanzen um $2''5$ zu klein gebe, so würde es gelöst sein. Allein ich kann durchaus nichts entdecken, was einen Fehler immer in einerlei Sinn hervorbrächte. Besonders will ich noch ausdrücklich bemerken, dass die Gesichtslinie des Fernrohrs mit der Ebene des Instruments auf das genaueste nach einem mir eigenthümlichen Verfahren parallel gestellt ist. Dass zugleich auf genaue Verticalität der Ebene gesehen ist, und dass bei der Berechnung nichts versäumt ist (z.\,B.\ sind allemal die vom Biquadrat abhängenden Glieder mit in Betracht gezogen), brauche ich nicht zu erinnern: allerdings macht die Vernachlässigung jeder dieser Vorsichtsregeln die Zenithdistanz zu klein. Meine Beobachtungen des letzten Solstitiums waren folgende:

\begin{center}
\fontsize{10}{12}\selectfont\setlength{\tabcolsep}{5pt}\renewcommand{\arraystretch}{1.15}\begin{tabular}{l l ccc l}
1814 Dec.\ 17 & W.\ Z.\ D.\ & $74^\circ$ & $53'$ & $48''53$ & 12 Beob. \\
19 & & 74 & 57 & 35.13 & 14 \\
21 & & 74 & 59 & 24.02 & 14 \\
22 & & 74 & 59 & 38.63 & 14 \\
23 & & 74 & 59 & 21.10 & 14 \\
24 & & 74 & 58 & 36.74 & 14 \\
25 & & 74 & 57 & 22.29 & 14 \\
29 & & 74 & 57 & 51.39 & 12 \\
\end{tabular}
\end{center}

\par\indent An jedem Tage waren die Beobachtungen in 2 oder 3 Reihen getheilt, die einzeln meistens kaum eine, nie zwei Secunden vom Mittel abwichen.

\par\indent Zur ganz scharfen Berechnung dieser Beobachtungen hat es mir noch an Zeit gefehlt. Vorläufig aber folgt daraus, die Z.\ D.\ im Solstitium $74^\circ\,59'\,36''6$, welches mit der Polhöhe $51^\circ\,31'\,55''0$ eine Schiefe gibt, die $5''$ kleiner ist als sie in unsern besten Tafeln angenommen wird. Die Polhöhe gründet sich auf mehr als 200 Beobachtungen des Polarsterns in der obern und untern Culmination. Letztere gab $51^\circ\,31'\,55''5$, erstere $51^\circ\,31'\,54''5$, wenn ich Hrn.\ \textsc{Bessel}'s Tafeln zum Grunde legte.

\Needspace{17\baselineskip}\begin{flushright}
Göttingen, den 24.\ April
\end{flushright}

\par\indent Den \textsc{Olbers}'schen Cometen werden Sie nun ohne Zweifel selbst längst aufgefunden und beobachtet haben. Ich fand ihn zwar gleich am 13.\ März, wo ich die erste Nachricht von seiner Entdeckung erhielt, auf; allein das äusserst ungünstige Wetter erlaubte nicht eher als am 20.\ eine wirkliche Beobachtung. Meine Beobachtungen aus dem März sind folgende:

\begin{center}\fontsize{9}{11}\selectfont\setlength{\tabcolsep}{3pt}\renewcommand{\arraystretch}{1.15}\begin{tabular}{l|ccc|ccc|ccc}
1815 & \multicolumn{3}{c|}{Mittlere Zeit} & \multicolumn{3}{c|}{\shortstack{Scheinb. ger.\\Aufsteigung}} & \multicolumn{3}{c}{\shortstack{Scheinb. Abw.\\Nordlich}} \\

\hline
März 20 & $10^{\mathrm u}$ & $33^{\mathrm m}$ & $6^{\mathrm s}$ & $54^\circ$ & $7'$ & $1''$ & $39^\circ$ & $7'$ & $47''$ \\
21 & 10 & 11 & 37 & 54 & 33 & 21 & 39 & 36 & 57 \\
21 & 10 & 57 & 0 & 54 & 34 & 21 & 39 & 38 & 8 \\
25 & 9 & 46 & 38 & 56 & 28 & 50 & 41 & 38 & 5 \\
30 & 9 & 50 & 57 & 59 & 14 & 3 & 44 & 10 & 27 \\
\end{tabular}\end{center}

\Needspace{20\baselineskip}\SourceLeaf{390}\par\indent Die zweite Beobachtung vom 21. ist von Hrn. \textsc{Enke}, welcher auch die sämmtlichen Vergleichungs-Sterne
aus der \emph{Histoire Céleste} reducirt hat. Meine Beobachtungen vom April sind noch nicht reducirt. ---
Den 31. März berechnete ich aus obigen Beobachtungen folgende vorläufige parabolische Elemente,
welche auch noch jetzt nicht viel abweichen werden.

\begin{center}\fontsize{9}{11}\selectfont $\text{Sonnennähe}\left\{\begin{array}{lcr}
\text{Zeit 1815. April 24.} & .\;.\;.\;. & 16^{\mathrm u}\;37^{\mathrm m}\;34^{\mathrm s}\;\text{Göttinger Zeit}\\
\text{Länge} & .\;.\;.\;.\;.\;.\;.\;.\;. & 146^\circ\;7\;2\\
\text{Abstand} & .\;.\;.\;.\;.\;.\;.\;. & 1.24738\\
\end{array}\right.$\end{center}
\begin{center}\fontsize{10}{12}\selectfont\setlength{\tabcolsep}{3pt}\renewcommand{\arraystretch}{1.15}\begin{tabular}{lcr}
Aufsteigender Knoten & $.\;.\;.\;.\;.\;.\;.\;.$ & $82\;43\;4$\\
Neigung der Bahn & $.\;.\;.\;.\;.\;.\;.\;.\;.$ & $45\;8\;55$\\
\end{tabular}\end{center}
\noindent Bewegung rechtläufig.

\par\indent Ich hoffe, dass dieser Comet noch bis zum Juli sich beobachten, und demnächst seine Bahn mit
vieler Zuverlässigkeit elliptisch berechnen lassen wird.

\par\indent Die Juno war um die Zeit ihrer Opposition nur 9. bis 10. Grösse; ich habe sie zweimal gut
beobachtet, aber die Beobachtungen noch nicht reducirt.

\par\indent Von unserm trefflichen Heliometer hat bisher, wegen Mangels einer festen Aufstellung, noch
wenig Gebrauch gemacht werden können. Vorgestern ist endlich das seit beinahe einem Jahre
bestellte, und äusserst schön und sinnreich von \textsc{Frauenhofer} ausgeführte parallatische Stativ
angekommen.
